% GENERATED FILE. Edit canonical lessons, then run npm run generate:books.
% canonical-source-hash: fnv1a64:b25416101290f10f
% canonical-lessons: ES-C617-decision, ES-C617-opcion, ES-C617-dilema, ES-C617-prioridad, ES-C617-criterio

\chapter{Elegir}
\label{ch:es-elegir}
\hlchaptermodality{617}

\begin{chapteropening}
\textbf{By the end of this chapter:} \emph{I can talk about a decision, the options and what comes first for me.}

Complete the last lesson of chapter 617: I can talk about a decision, the options and what comes first for me.
\end{chapteropening}

\section[la decisión]{la decisión — the decision}

\label{lesson:ES-C617-decision}



\begin{quote}
Before the new one: say the Spanish for the setback, then the Spanish for the alternative.
\end{quote}

\subsection*{You'll want to know: la decisión}
\textbf{la decisión} — \textquotedblleft{}the decision\textquotedblright{}.

\begin{quote}
Es una decisión difícil.
\end{quote}

\begin{quote}
\emph{It is a difficult decision.}
\end{quote}

\emph{Tomar una decisión} — to make a decision.

\begin{cousinweb}
\emph{decisión} is a learned borrowing from Latin \emph{decisio} 'a settlement', built on \emph{decidere} 'to cut off, to settle', from \emph{caedere} 'to cut' (EtymDB).
\end{cousinweb}

\begin{grammarlens}[title={What you've built}]
One more word for telling, explaining and describing at B1.
\end{grammarlens}

\subsection*{Your turn}
\begin{itemize}
\raggedright
  \item \emph{Say it:} \emph{la decisión}
  \item \emph{Say it:} \emph{la decisión}, once more
  \item \emph{Say it:} a sentence of your own with \emph{la decisión}
  \item \emph{Recall:} say the Spanish for the setback, then the Spanish for the alternative, then say \emph{la decisión} again
\end{itemize}

\begin{tcolorbox}[breakable,colback=teal!4,colframe=teal!35!black,title={Before you move on}]
What does la decisión mean? (\textquotedblleft{}The decision\textquotedblright{}.) Say it once more.
\end{tcolorbox}
\section[la opción]{la opción — the option, the choice}

\label{lesson:ES-C617-opcion}



\begin{quote}
Before the new one: say the Spanish for the alternative, then the Spanish for the decision.
\end{quote}

\subsection*{You'll want to know: la opción}
\textbf{la opción} — \textquotedblleft{}the option, the choice\textquotedblright{}.

\begin{quote}
Tenemos dos opciones.
\end{quote}

\begin{quote}
\emph{We have two options.}
\end{quote}

\emph{No tener opción} — to have no choice.

\begin{cousinweb}
\emph{opción} is a learned borrowing from Latin \emph{optio} 'free choice', built on \emph{optare} 'to choose, to wish' (Lewis \& Short).
\end{cousinweb}

\begin{grammarlens}[title={What you've built}]
One more word for telling, explaining and describing at B1.
\end{grammarlens}

\subsection*{Your turn}
\begin{itemize}
\raggedright
  \item \emph{Say it:} \emph{la opción}
  \item \emph{Say it:} \emph{la opción}, once more
  \item \emph{Say it:} a sentence of your own with \emph{la opción}
  \item \emph{Recall:} say the Spanish for the alternative, then the Spanish for the decision, then say \emph{la opción} again
\end{itemize}

\begin{tcolorbox}[breakable,colback=teal!4,colframe=teal!35!black,title={Before you move on}]
What does la opción mean? (\textquotedblleft{}The option, the choice\textquotedblright{}.) Say it once more.
\end{tcolorbox}
\section[el dilema]{el dilema — the dilemma}

\label{lesson:ES-C617-dilema}



\begin{quote}
Before the new one: say the Spanish for the decision, then the Spanish for the option.
\end{quote}

\subsection*{You'll want to know: el dilema}
\textbf{el dilema} — \textquotedblleft{}the dilemma\textquotedblright{}.

\begin{quote}
Tengo un dilema: trabajar o estudiar.
\end{quote}

\begin{quote}
\emph{I have a dilemma: to work or to study.}
\end{quote}

\emph{Estar ante un dilema} — to face a dilemma.

\begin{cousinweb}
\emph{dilema} comes from Latin \emph{dilemma}, taken from Greek \emph{dilēmma} 'a double premise', from \emph{dis} 'twice' and \emph{lēmma} 'premise' (EtymDB).
\end{cousinweb}

\begin{grammarlens}[title={What you've built}]
One more word for telling, explaining and describing at B1.
\end{grammarlens}

\subsection*{Your turn}
\begin{itemize}
\raggedright
  \item \emph{Say it:} \emph{el dilema}
  \item \emph{Say it:} \emph{el dilema}, once more
  \item \emph{Say it:} a sentence of your own with \emph{el dilema}
  \item \emph{Recall:} say the Spanish for the decision, then the Spanish for the option, then say \emph{el dilema} again
\end{itemize}

\begin{tcolorbox}[breakable,colback=teal!4,colframe=teal!35!black,title={Before you move on}]
What does el dilema mean? (\textquotedblleft{}The dilemma\textquotedblright{}.) Say it once more.
\end{tcolorbox}
\section[la prioridad]{la prioridad — the priority}

\label{lesson:ES-C617-prioridad}



\begin{quote}
Before the new one: say the Spanish for the option, then the Spanish for the dilemma.
\end{quote}

\subsection*{You'll want to know: la prioridad}
\textbf{la prioridad} — \textquotedblleft{}the priority\textquotedblright{}.

\begin{quote}
Mi familia es mi prioridad.
\end{quote}

\begin{quote}
\emph{My family is my priority.}
\end{quote}

\emph{Dar prioridad a} — to put first.

\begin{cousinweb}
\emph{prioridad} is a learned word usually traced to Medieval Latin \emph{prioritas}, built on Latin \emph{prior} 'earlier, former' (Lewis \& Short).
\end{cousinweb}

\begin{grammarlens}[title={What you've built}]
One more word for telling, explaining and describing at B1.
\end{grammarlens}

\subsection*{Your turn}
\begin{itemize}
\raggedright
  \item \emph{Say it:} \emph{la prioridad}
  \item \emph{Say it:} \emph{la prioridad}, once more
  \item \emph{Say it:} a sentence of your own with \emph{la prioridad}
  \item \emph{Recall:} say the Spanish for the option, then the Spanish for the dilemma, then say \emph{la prioridad} again
\end{itemize}

\begin{tcolorbox}[breakable,colback=teal!4,colframe=teal!35!black,title={Before you move on}]
What does la prioridad mean? (\textquotedblleft{}The priority\textquotedblright{}.) Say it once more.
\end{tcolorbox}
\section[el criterio]{el criterio — the criterion; judgement}

\label{lesson:ES-C617-criterio}



\begin{quote}
Before the new one: say the Spanish for the dilemma, then the Spanish for the priority.
\end{quote}

\subsection*{You'll want to know: el criterio}
\textbf{el criterio} — \textquotedblleft{}the criterion; judgement\textquotedblright{}.

\begin{quote}
¿Con qué criterio eliges?
\end{quote}

\begin{quote}
\emph{On what basis do you choose?}
\end{quote}

\emph{Tener criterio} — to have good judgement.

\begin{cousinweb}
\emph{criterio} is a learned borrowing from Greek \emph{kritērion} 'a test, a means for judging', built on \emph{krinein} 'to judge'; EtymDB does not list it, but it gives the same root for \emph{crítico}.
\end{cousinweb}

\begin{grammarlens}[title={What you've built}]
One more word for telling, explaining and describing at B1.
\end{grammarlens}

\subsection*{Your turn}
\begin{itemize}
\raggedright
  \item \emph{Say it:} \emph{el criterio}
  \item \emph{Say it:} \emph{el criterio}, once more
  \item \emph{Say it:} a sentence of your own with \emph{el criterio}
  \item \emph{Recall:} say the Spanish for the dilemma, then the Spanish for the priority, then say \emph{el criterio} again
\end{itemize}

\begin{tcolorbox}[breakable,colback=teal!4,colframe=teal!35!black,title={Before you move on}]
What does el criterio mean? (\textquotedblleft{}The criterion; judgement\textquotedblright{}.) Say it once more.
\end{tcolorbox}
